% Part: sets-functions-relations
% Chapter: size-of-sets-complete

\documentclass[../../../include/open-logic-chapter]{subfiles}

\begin{document}

\olchapter{sfr}{siz}{حجم المجموعات}

\begin{editorial}
يناقش هذا الفصل التعدادات، وقابلية العد، وعدم قابلية العد. وترد عدة
أقسام في نسختين: نسخة أيسر وأقرب إلى المستوى التمهيدي، تعدّ التعدادات
قوائم أو دوال شاملة من $\PosInt$؛ ونسخة أكثر تجريدًا تعرّف التعدادات
بأنها دوال تقابلية مع $\Nat$.
\end{editorial}

\olimport{introduction}

\olimport{enumerability}

\olimport{zig-zag}

\olimport{pairing}

\olimport{pairing-alt}

\olimport{non-enumerability}

\olimport{reduction}

\olimport{equinumerous-sets}

\olimport{comparing-size}

\olimport{schroder-bernstein}

\begin{editorial}
الأقسام الآتية: \olref[sfr][siz][enm-alt]{sec}،
و\olref[sfr][siz][nen-alt]{sec}، و\olref[sfr][siz][red-alt]{sec}، هي
نسخ بديلة من \olref[sfr][siz][enm]{sec}،
و\olref[sfr][siz][nen]{sec}، و\olref[sfr][siz][red]{sec}، أعدّها تيم باتن
لاستعمالها في كتابه «نظرية المجموعات المفتوحة». وهي أكثر تقدمًا
بقليل، وتستعمل تعريفًا مختلفًا لقابلية التعداد، أنسب لسياق نظرية
المجموعات (أي وجود دالة تقابلية مع $\Nat$ أو مع مقطع ابتدائي، بدلًا من
قابلية السرد أو كون المجموعة مدى دالة شاملة من $\PosInt$).
\end{editorial}

\olimport{enumerability-alt}
\olimport{non-enumerability-alt}
\olimport{reduction-alt}

\OLEndChapterHook

\end{document}
